\section{Conclusion}


Bridge security reduces to message layer security. The authenticated cross-chain message protocol eliminates the three root causes of bridge exploits: weak authentication (replaced by MPC threshold attestation), missing replay protection (replaced by five independent mechanisms), and absent backing verification (replaced by monotonic backing attestation with auto-pause). The protocol is implemented as the Lux Teleport bridge, supporting 270+ chains with zero forged, replayed, or incorrectly processed messages in production.

\begin{thebibliography}{9}
\bibitem{cggmp2021} Canetti, Gennaro, Goldfeder, Makriyannis, Peled, ``UC Non-Interactive, Proactive, Threshold ECDSA with Identifiable Aborts,'' CCS 2020.
\bibitem{frost2020} Komlo, Goldberg, ``FROST: Flexible Round-Optimized Schnorr Threshold Signatures,'' SAC 2020.
\bibitem{lux-teleport} Lux Industries, Inc., ``Teleport Universal Bridge Architecture,'' 2025.
\bibitem{lux-warp} Lux Industries, Inc., ``Warp Messaging: Authenticated Cross-Chain Communication,'' 2025.
\bibitem{lux-triple} Lux Industries, Inc., ``Triple-Proof Quantum Finality,'' 2026.
\bibitem{ronin} US Treasury, ``OFAC Designates Lazarus Group Ethereum Address in Connection with Ronin Bridge Exploit,'' 2022.
\end{thebibliography}

